\section{Quad X PWM Mixer}

Mixer là tầng nhận các lệnh PWM trừu tượng và tạo PWM cụ thể cho từng motor.
Cấu hình của drone là quad X với thứ tự motor:

\begin{center}
\begin{tabular}{|c|c|c|}
\hline
Motor & Vị trí trên frame & Tọa độ tương đối \\
\hline
1 & front-left & $(+arm,\ +arm,\ z_m)$ \\
\hline
2 & front-right & $(+arm,\ -arm,\ z_m)$ \\
\hline
3 & rear-right & $(-arm,\ -arm,\ z_m)$ \\
\hline
4 & rear-left & $(-arm,\ +arm,\ z_m)$ \\
\hline
\end{tabular}
\end{center}

Trong code:

\begin{equation}
    arm = 0.159,\quad z_m = 0.04
\end{equation}

\subsection{Công thức mixer}

Mixer nhận:

\begin{equation}
    PWM_{base},\quad
    \Delta PWM_\phi,\quad
    \Delta PWM_\theta,\quad
    \Delta PWM_\psi
\end{equation}

và tính:

\begin{equation}
    PWM_1
    =
    PWM_{base}
    +
    \Delta PWM_\phi
    -
    \Delta PWM_\theta
    -
    \Delta PWM_\psi
\end{equation}

\begin{equation}
    PWM_2
    =
    PWM_{base}
    -
    \Delta PWM_\phi
    -
    \Delta PWM_\theta
    +
    \Delta PWM_\psi
\end{equation}

\begin{equation}
    PWM_3
    =
    PWM_{base}
    -
    \Delta PWM_\phi
    +
    \Delta PWM_\theta
    -
    \Delta PWM_\psi
\end{equation}

\begin{equation}
    PWM_4
    =
    PWM_{base}
    +
    \Delta PWM_\phi
    +
    \Delta PWM_\theta
    +
    \Delta PWM_\psi
\end{equation}

Có thể được viết dưới dạng ma trận:

\begin{equation}
    \begin{bmatrix}
        PWM_1 \\
        PWM_2 \\
        PWM_3 \\
        PWM_4
    \end{bmatrix}
    = PWM_{base}
    \begin{bmatrix}
        1 \\
        1 \\
        1 \\
        1
    \end{bmatrix}
    + \Delta PWM_{\phi}
    \begin{bmatrix}
        +1 \\
        -1 \\
        -1 \\
        +1
    \end{bmatrix}
    + \Delta PWM_{\theta}
    \begin{bmatrix}
        -1 \\
        -1 \\
        +1 \\
        +1
    \end{bmatrix}
    + \Delta PWM_{\psi}
    \begin{bmatrix}
        -1 \\
        +1 \\
        -1 \\
        +1
    \end{bmatrix}
\end{equation}
Sau đó từng PWM được giới hạn:

\begin{equation}
    1100 \leq PWM_i \leq 1940
\end{equation}

\subsection{Ý nghĩa dấu trong mixer}

Các dấu cộng/trừ trong mixer không được chọn ngẫu nhiên. Chúng xuất phát từ
vị trí của từng motor trên thân drone và chiều quay của từng cánh quạt.

Ý tưởng chính của mixer là:

\begin{equation}
    PWM_{base}
    \rightarrow
    \text{tạo lực nâng chung}
\end{equation}

\begin{equation}
    \Delta PWM_\phi,\ \Delta PWM_\theta,\ \Delta PWM_\psi
    \rightarrow
    \text{tạo chênh lệch lực để sinh moment quay}
\end{equation}

Trong đó:
\begin{itemize}
    \item $\Delta PWM_\phi$: hiệu chỉnh roll quanh trục $x$.
    \item $\Delta PWM_\theta$: hiệu chỉnh pitch quanh trục $y$.
    \item $\Delta PWM_\psi$: hiệu chỉnh yaw quanh trục $z$.
\end{itemize}

\subsubsection{Quy ước hệ trục và vị trí motor}

Trong project này, hệ trục body của drone được quy ước:

\begin{equation}
    +x:\ \text{phía trước},\qquad
    +y:\ \text{phía trái},\qquad
    +z:\ \text{hướng lên}
\end{equation}

Bốn motor được đặt tại các vị trí:

\begin{equation}
    M_1: (+x,+y)
\end{equation}

\begin{equation}
    M_2: (+x,-y)
\end{equation}

\begin{equation}
    M_3: (-x,-y)
\end{equation}

\begin{equation}
    M_4: (-x,+y)
\end{equation}

Do đó:
\begin{itemize}
    \item $M_1$ và $M_4$ nằm bên trái drone vì có $y>0$.
    \item $M_2$ và $M_3$ nằm bên phải drone vì có $y<0$.
    \item $M_1$ và $M_2$ nằm phía trước vì có $x>0$.
    \item $M_3$ và $M_4$ nằm phía sau vì có $x<0$.
\end{itemize}

\subsubsection{Vì sao chênh lệch lực tạo ra moment?}

Mỗi motor tạo ra một lực nâng $F_i$ theo hướng lên. Vì motor không nằm ở tâm
drone, lực nâng này sẽ tạo moment quay quanh tâm drone.

Công thức moment tổng quát là:

\begin{equation}
    \boldsymbol{\tau}
    =
    \mathbf{r}
    \times
    \mathbf{F}
\end{equation}

Trong đó:
\begin{itemize}
    \item $\mathbf{r}$ là vector từ tâm drone đến vị trí motor.
    \item $\mathbf{F}$ là lực nâng do motor tạo ra.
    \item $\boldsymbol{\tau}$ là moment làm drone quay.
\end{itemize}

Nếu lực đặt đúng tại tâm drone thì $\mathbf{r}=0$, do đó không tạo moment.
Ngược lại, nếu lực đặt lệch khỏi tâm, tức là $\mathbf{r}\neq 0$, lực đó sẽ làm
drone quay.

Với motor thứ $i$, ta có thể viết:

\begin{equation}
    \mathbf{r}_i
    =
    \begin{bmatrix}
        x_i\\
        y_i\\
        0
    \end{bmatrix},
    \qquad
    \mathbf{F}_i
    =
    \begin{bmatrix}
        0\\
        0\\
        F_i
    \end{bmatrix}
\end{equation}

Moment do motor đó tạo ra là:

\begin{equation}
    \boldsymbol{\tau}_i
    =
    \mathbf{r}_i
    \times
    \mathbf{F}_i
    =
    \begin{bmatrix}
        y_iF_i\\
        -x_iF_i\\
        0
    \end{bmatrix}
\end{equation}

Do đó:

\begin{equation}
    \tau_x = y_iF_i
\end{equation}

\begin{equation}
    \tau_y = -x_iF_i
\end{equation}

Nghĩa là:
\begin{itemize}
    \item Moment roll quanh trục $x$ phụ thuộc vào $y_iF_i$.
    \item Moment pitch quanh trục $y$ phụ thuộc vào $-x_iF_i$.
\end{itemize}

\subsubsection{Ý nghĩa của $PWM_{base}$}

$PWM_{base}$ là PWM nền được cộng đều vào cả bốn motor:

\begin{equation}
    PWM_1 = PWM_2 = PWM_3 = PWM_4 = PWM_{base}
\end{equation}

Khi bốn motor nhận PWM gần như bằng nhau, chúng tạo lực nâng gần như bằng nhau.
Khi đó các moment roll và pitch triệt tiêu nhau, drone chủ yếu hover, bay lên
hoặc bay xuống.

Có thể hiểu:

\begin{equation}
    PWM_{base}
    \rightarrow
    \text{điều khiển tổng lực nâng}
\end{equation}

\subsubsection{Ý nghĩa chung của các delta PWM}

Các thành phần delta PWM không nhằm tăng đều cả bốn motor, mà nhằm tăng motor
này và giảm motor khác để tạo mất cân bằng lực.

Dạng tổng quát của mixer là:

\begin{equation}
\begin{bmatrix}
    PWM_1\\
    PWM_2\\
    PWM_3\\
    PWM_4
\end{bmatrix}
=
PWM_{base}
\begin{bmatrix}
    1\\
    1\\
    1\\
    1
\end{bmatrix}
+
\Delta PWM_\phi
\begin{bmatrix}
    +1\\
    -1\\
    -1\\
    +1
\end{bmatrix}
+
\Delta PWM_\theta
\begin{bmatrix}
    -1\\
    -1\\
    +1\\
    +1
\end{bmatrix}
+
\Delta PWM_\psi
\begin{bmatrix}
    -1\\
    +1\\
    -1\\
    +1
\end{bmatrix}
\end{equation}

Các vector dấu của delta PWM đều có tổng bằng $0$:

\begin{equation}
    (+1)+(-1)+(-1)+(+1)=0
\end{equation}

\begin{equation}
    (-1)+(-1)+(+1)+(+1)=0
\end{equation}

\begin{equation}
    (-1)+(+1)+(-1)+(+1)=0
\end{equation}

Vì vậy, về ý tưởng, các delta PWM không làm thay đổi nhiều tổng lực nâng. Chúng
chủ yếu tạo chênh lệch lực giữa các motor để tạo moment quay.

\subsubsection{Roll correction tạo moment roll}

Roll là chuyển động quay quanh trục $x$. Moment roll tổng do bốn motor tạo ra
có thể viết là:

\begin{equation}
    \tau_x
    =
    \sum_{i=1}^{4} y_iF_i
\end{equation}

Trong đó $y_i$ là vị trí của motor theo trục $y$, còn $F_i$ là lực nâng của
motor đó.

Với layout hiện tại:

\begin{equation}
    y_1 = +l,\qquad y_4 = +l
\end{equation}

\begin{equation}
    y_2 = -l,\qquad y_3 = -l
\end{equation}

Nếu $\Delta PWM_\phi > 0$, mixer chọn:

\begin{equation}
    M_1:\ +,\qquad
    M_2:\ -,\qquad
    M_3:\ -,\qquad
    M_4:\ +
\end{equation}

Nghĩa là:
\begin{itemize}
    \item $M_1$ và $M_4$ tăng lực nâng.
    \item $M_2$ và $M_3$ giảm lực nâng.
\end{itemize}

Ta có thể viết gần đúng:

\begin{equation}
    F_1 = F_0 + \Delta F
\end{equation}

\begin{equation}
    F_4 = F_0 + \Delta F
\end{equation}

\begin{equation}
    F_2 = F_0 - \Delta F
\end{equation}

\begin{equation}
    F_3 = F_0 - \Delta F
\end{equation}

Khi đó:

\begin{equation}
\begin{aligned}
    \tau_x
    &=
    y_1F_1 + y_2F_2 + y_3F_3 + y_4F_4\\
    &=
    l(F_0+\Delta F)
    -l(F_0-\Delta F)
    -l(F_0-\Delta F)
    +l(F_0+\Delta F)
\end{aligned}
\end{equation}

Các thành phần $F_0$ triệt tiêu nhau, còn lại:

\begin{equation}
    \tau_x = 4l\Delta F
\end{equation}

Vì vậy, nếu $\Delta F>0$ thì:

\begin{equation}
    \tau_x > 0
\end{equation}

Nghĩa là chênh lệch lực trái--phải tạo ra moment roll dương.

\subsubsection{Vì sao roll correction không làm đổi nhiều tổng lực nâng?}

Tổng lực nâng của bốn motor là:

\begin{equation}
    F_{total}
    =
    F_1 + F_2 + F_3 + F_4
\end{equation}

Thay các lực sau khi roll correction:

\begin{equation}
\begin{aligned}
    F_{total}
    &=
    (F_0+\Delta F)
    +(F_0-\Delta F)
    +(F_0-\Delta F)
    +(F_0+\Delta F)
\end{aligned}
\end{equation}

Các thành phần $\Delta F$ triệt tiêu:

\begin{equation}
    +\Delta F-\Delta F-\Delta F+\Delta F=0
\end{equation}

Do đó:

\begin{equation}
    F_{total}=4F_0
\end{equation}

Vì vậy roll correction gần như không làm đổi tổng lực nâng. Nó chủ yếu làm lực
bên trái lớn hơn bên phải, từ đó tạo moment roll.

\subsubsection{Pitch correction tạo moment pitch}

Pitch là chuyển động quay quanh trục $y$. Moment pitch tổng có dạng:

\begin{equation}
    \tau_y
    =
    \sum_{i=1}^{4} -x_iF_i
\end{equation}

Motor phía trước có $x>0$ là $M_1$ và $M_2$. Motor phía sau có $x<0$ là $M_3$
và $M_4$.

Nếu $\Delta PWM_\theta > 0$, mixer chọn:

\begin{equation}
    M_1:\ -,\qquad
    M_2:\ -,\qquad
    M_3:\ +,\qquad
    M_4:\ +
\end{equation}

Nghĩa là:
\begin{itemize}
    \item $M_1$ và $M_2$ giảm lực nâng.
    \item $M_3$ và $M_4$ tăng lực nâng.
\end{itemize}

Khi đó phía sau đẩy mạnh hơn phía trước. Vì phía sau có $x<0$, theo công thức:

\begin{equation}
    \tau_y = \sum_{i=1}^{4} -x_iF_i
\end{equation}

sự tăng lực ở phía sau sẽ tạo moment pitch dương theo quy ước của project.

Do đó vector dấu của pitch correction là:

\begin{equation}
    \Delta PWM_\theta
    \begin{bmatrix}
        -1\\
        -1\\
        +1\\
        +1
    \end{bmatrix}
\end{equation}

\subsubsection{Yaw correction tạo moment yaw}

Yaw là chuyển động quay quanh trục $z$. Phần yaw khó hơn roll và pitch vì yaw
không được tạo chủ yếu bằng cách tạo chênh lệch lực nâng trái--phải hoặc
trước--sau. Yaw chủ yếu được tạo bởi torque phản lực của cánh quạt.

Khi một motor làm cánh quạt quay, motor tác dụng một torque lên cánh quạt.
Theo định luật III Newton, cánh quạt cũng tác dụng ngược lại một torque lên
thân drone. Vì vậy, khi cánh quạt quay một chiều, thân drone có xu hướng quay chiều ngược lại. 

Đây là nguồn gốc của yaw torque.

Trong project, chiều quay motor được quy ước bởi vector:

\begin{equation}
    d =
    \begin{bmatrix}
        +1\\
        -1\\
        +1\\
        -1
    \end{bmatrix}
\end{equation}

Nghĩa là $M_1$ và $M_3$ quay cùng một chiều, còn $M_2$ và $M_4$ quay chiều
ngược lại.

Gọi $Q_i$ là reaction torque của motor thứ $i$. Torque yaw tổng tác dụng lên
body được mô hình hóa bởi:

\begin{equation}
    M_z
    =
    \sum_{i=1}^{4} -d_iQ_i
\end{equation}

Viết tường minh:

\begin{equation}
    M_z
    =
    -d_1Q_1
    -d_2Q_2
    -d_3Q_3
    -d_4Q_4
\end{equation}

Vì:

\begin{equation}
    d_1=+1,\quad
    d_2=-1,\quad
    d_3=+1,\quad
    d_4=-1
\end{equation}

nên:

\begin{equation}
\begin{aligned}
    M_z
    &=
    -(+1)Q_1
    -(-1)Q_2
    -(+1)Q_3
    -(-1)Q_4\\
    &=
    -Q_1 + Q_2 - Q_3 + Q_4
\end{aligned}
\end{equation}

Khi hover đều, bốn motor tạo reaction torque gần như bằng nhau:

\begin{equation}
    Q_1 = Q_2 = Q_3 = Q_4 = Q_0
\end{equation}

Thay vào công thức trên:

\begin{equation}
\begin{aligned}
    M_z
    &=
    -Q_0 + Q_0 - Q_0 + Q_0\\
    &=
    0
\end{aligned}
\end{equation}

Vì vậy, khi bốn motor chạy đều, các reaction torque quanh trục $z$ triệt tiêu
nhau và drone không tự quay yaw.

\textbf{Tạo yaw moment dương}

Từ công thức:

\begin{equation}
    M_z = -Q_1 + Q_2 - Q_3 + Q_4
\end{equation}

ta thấy $Q_2$ và $Q_4$ đóng góp dấu dương vào $M_z$, còn $Q_1$ và $Q_3$ đóng
góp dấu âm vào $M_z$.

Do đó, để tạo yaw moment dương:

\begin{equation}
    M_z > 0
\end{equation}

ta cần tăng reaction torque của nhóm $M_2$, $M_4$ và giảm reaction torque của
nhóm $M_1$, $M_3$:

\begin{equation}
    Q_2 \uparrow,\quad Q_4 \uparrow
\end{equation}

\begin{equation}
    Q_1 \downarrow,\quad Q_3 \downarrow
\end{equation}

Vì reaction torque $Q_i$ tăng khi PWM motor tăng, điều này tương ứng với:

\begin{equation}
    PWM_2 \uparrow,\quad PWM_4 \uparrow
\end{equation}

\begin{equation}
    PWM_1 \downarrow,\quad PWM_3 \downarrow
\end{equation}

Do đó, nếu $\Delta PWM_\psi > 0$, mixer chọn:

\begin{equation}
    M_1:\ -,\qquad
    M_2:\ +,\qquad
    M_3:\ -,\qquad
    M_4:\ +
\end{equation}

Hay viết dưới dạng vector dấu:

\begin{equation}
    \Delta PWM_\psi
    \begin{bmatrix}
        -1\\
        +1\\
        -1\\
        +1
    \end{bmatrix}
\end{equation}

\paragraph{Chứng minh bằng công thức}

Giả sử ban đầu reaction torque của bốn motor bằng nhau:

\begin{equation}
    Q_1 = Q_2 = Q_3 = Q_4 = Q_0
\end{equation}

Sau khi áp dụng yaw correction dương:

\begin{equation}
    Q_1 = Q_0 - \Delta Q
\end{equation}

\begin{equation}
    Q_2 = Q_0 + \Delta Q
\end{equation}

\begin{equation}
    Q_3 = Q_0 - \Delta Q
\end{equation}

\begin{equation}
    Q_4 = Q_0 + \Delta Q
\end{equation}

Thay vào công thức:

\begin{equation}
    M_z = -Q_1 + Q_2 - Q_3 + Q_4
\end{equation}

ta được:

\begin{equation}
\begin{aligned}
    M_z
    &=
    -(Q_0-\Delta Q)
    +(Q_0+\Delta Q)
    -(Q_0-\Delta Q)
    +(Q_0+\Delta Q)
\end{aligned}
\end{equation}

Khai triển:

\begin{equation}
\begin{aligned}
    M_z
    &=
    -Q_0+\Delta Q
    +Q_0+\Delta Q
    -Q_0+\Delta Q
    +Q_0+\Delta Q
\end{aligned}
\end{equation}

Các thành phần $Q_0$ triệt tiêu nhau:

\begin{equation}
    -Q_0 + Q_0 - Q_0 + Q_0 = 0
\end{equation}

Còn lại:

\begin{equation}
    M_z = 4\Delta Q
\end{equation}

Vì vậy, nếu:

\begin{equation}
    \Delta Q > 0
\end{equation}

thì:

\begin{equation}
    M_z > 0
\end{equation}

Nghĩa là drone tạo được yaw moment dương.

\paragraph{Vì sao yaw correction không làm đổi nhiều tổng lực nâng?}

Vector dấu của yaw correction là:

\begin{equation}
    \begin{bmatrix}
        -1\\
        +1\\
        -1\\
        +1
    \end{bmatrix}
\end{equation}

Tổng các phần tử của vector này bằng $0$:

\begin{equation}
    (-1)+(+1)+(-1)+(+1)=0
\end{equation}

Điều này có nghĩa là mixer giảm hai motor và tăng hai motor. Vì vậy, về ý
tưởng, tổng lực nâng không thay đổi nhiều. Tuy nhiên, do hai nhóm motor quay
ngược chiều nhau, reaction torque không còn cân bằng nữa, nên drone sẽ quay
quanh trục yaw.

\paragraph{So sánh yaw với roll và pitch}

Roll và pitch được tạo chủ yếu bởi chênh lệch lực nâng theo vị trí motor:

\begin{equation}
    \tau_x = \sum_{i=1}^{4} y_iF_i
\end{equation}

\begin{equation}
    \tau_y = \sum_{i=1}^{4} -x_iF_i
\end{equation}

Nghĩa là roll và pitch phụ thuộc vào motor nằm bên trái--phải hoặc trước--sau.

Yaw thì khác. Yaw chủ yếu phụ thuộc vào chiều quay của cánh quạt:

\begin{equation}
    M_z = \sum_{i=1}^{4} -d_iQ_i
\end{equation}

Do đó, yaw correction không chia motor theo trái--phải hoặc trước--sau, mà chia
theo hai nhóm motor quay ngược chiều:

\begin{equation}
    M_1,\ M_3
\end{equation}

và:

\begin{equation}
    M_2,\ M_4
\end{equation}

\paragraph{Cách nhớ nhanh}

Trong project:

\begin{equation}
    d =
    \begin{bmatrix}
        +1\\
        -1\\
        +1\\
        -1
    \end{bmatrix}
\end{equation}

và:

\begin{equation}
    M_z = \sum_{i=1}^{4} -d_iQ_i
\end{equation}

Do có dấu trừ trước $d_i$, ta có:

\begin{equation}
    d_i = +1
    \Rightarrow
    \text{đóng góp yaw âm}
\end{equation}

\begin{equation}
    d_i = -1
    \Rightarrow
    \text{đóng góp yaw dương}
\end{equation}

Vì vậy, muốn tạo yaw moment dương thì giảm nhóm $d_i=+1$ và tăng nhóm
$d_i=-1$:

\begin{equation}
    M_1:\ -,\qquad
    M_2:\ +,\qquad
    M_3:\ -,\qquad
    M_4:\ +
\end{equation}

Do đó:

\begin{equation}
    \Delta PWM_\psi
    \begin{bmatrix}
        -1\\
        +1\\
        -1\\
        +1
    \end{bmatrix}
\end{equation}

\paragraph{Tóm tắt}

Motor quay cánh quạt sẽ tạo reaction torque lên thân drone. Hai motor $M_1$,
$M_3$ tạo yaw torque một chiều, còn hai motor $M_2$, $M_4$ tạo yaw torque chiều
ngược lại. Khi hover đều, hai chiều torque này triệt tiêu nhau. Muốn drone quay
yaw, mixer tăng PWM của một nhóm và giảm PWM của nhóm còn lại. Nhờ vậy tổng lực
nâng gần như giữ nguyên, nhưng torque yaw không còn cân bằng, nên drone quay
quanh trục $z$.

\section{Motor Model}

Motor model mô phỏng đơn giản ESC, motor BLDC và propeller. Input là PWM, output
là lực đẩy, dòng điện, RPM và torque phản lực.

\subsection{Đổi PWM sang throttle}

PWM được đưa về tỉ lệ ga $u$:

\begin{equation}
    u
    =
    \frac{PWM-PWM_{min}}{PWM_{max}-PWM_{min}}
\end{equation}

Với:

\begin{equation}
    0 \leq u \leq 1
\end{equation}

Điện áp hiệu dụng đưa vào motor:

\begin{equation}
    V_m = uV_{bat}
\end{equation}

Trong code:

\begin{equation}
    V_{bat}=14.8\ \mathrm{V}
\end{equation}

\subsection{Dòng điện và torque motor}

Hằng số back-EMF được tính từ KV:

\begin{equation}
    K_e = \frac{60}{2\pi KV}
\end{equation}

Code dùng:

\begin{equation}
    K_t = K_e
\end{equation}

Dòng điện motor:

\begin{equation}
    I
    =
    \frac{V_m-K_e\omega}{R_m}
\end{equation}

Dòng điện được giới hạn:

\begin{equation}
    0 \leq I \leq I_{max}
\end{equation}

Torque motor:

\begin{equation}
    \tau_m
    =
    K_t \max(I-I_0,0)
\end{equation}

Trong đó $I_0$ là dòng không tải.

\subsection{Động lực học rotor}

Torque tải khí động của propeller:

\begin{equation}
    \tau_{load}
    =
    k_q\omega^2
\end{equation}

Gia tốc góc rotor:

\begin{equation}
    \dot{\omega}
    =
    \frac{\tau_m-\tau_{load}}{J}
\end{equation}

Cập nhật rời rạc:

\begin{equation}
    \omega(k)
    =
    \omega(k-1)
    +
    \dot{\omega}\Delta t
\end{equation}

Sau đó $\omega$ được kẹp trong khoảng:

\begin{equation}
    0 \leq \omega \leq \omega_{max}
\end{equation}

Trong code:

\begin{equation}
    \omega_{max}=1800\ \mathrm{rad/s}
\end{equation}

\subsection{Lực đẩy, torque phản lực và RPM}

Lực đẩy propeller:

\begin{equation}
    T = k_f\omega^2
\end{equation}

Torque phản lực:

\begin{equation}
    Q = k_q\omega^2
\end{equation}

RPM:

\begin{equation}
    RPM
    =
    \frac{\omega 60}{2\pi}
\end{equation}

Các tham số mặc định của motor model:

\begin{center}
\begin{tabular}{|l|l|}
\hline
Tham số & Giá trị \\
\hline
$KV$ & $1100$ \\
\hline
$R_m$ & $0.073\ \Omega$ \\
\hline
$I_0$ & $0.9\ \mathrm{A}$ \\
\hline
$I_{max}$ & $30\ \mathrm{A}$ \\
\hline
$J$ & $2.5\times10^{-5}$ \\
\hline
$k_f$ & $4.0\times10^{-6}$ \\
\hline
$k_q$ & $1.0\times10^{-7}$ \\
\hline
\end{tabular}
\end{center}

\section{Apply lực vào Isaac Sim}

Sau khi motor model trả về lực đẩy $T_i$, controller apply lực lên rigid body
\texttt{base\_link} tại vị trí của từng motor:

\begin{equation}
    \mathbf{F}_i
    =
    \begin{bmatrix}
        0 \\
        0 \\
        T_i
    \end{bmatrix}
\end{equation}

Vì lực được đặt lệch khỏi tâm drone, Isaac Sim tự sinh moment roll/pitch tương
ứng. Đây là lý do mixer không cần apply moment roll/pitch trực tiếp; nó chỉ cần
tạo chênh lệch lực motor.

\subsection{Torque phản lực yaw}

Mỗi propeller tạo torque phản lực $Q_i$. Với chiều quay motor:

\begin{equation}
    d = [1,\ -1,\ 1,\ -1]
\end{equation}

Torque yaw tổng lên body được tính:

\begin{equation}
    M_z
    =
    \sum_{i=1}^{4} -d_iQ_i
\end{equation}

Trong code, torque yaw này được đổi thành một cặp lực ngang đặt tại hai điểm
đối xứng:

\begin{equation}
    F_y
    =
    \frac{M_z}{2arm}
\end{equation}

Một lực $+F_y$ đặt tại $x=+arm$, lực còn lại $-F_y$ đặt tại $x=-arm$. Hai lực
này triệt tiêu tổng lực tuyến tính nhưng tạo moment quanh trục $z$.

\subsection{Kết nối giữa motor model và simulator}

Mối quan hệ cuối cùng giữa PWM và chuyển động drone là:

\begin{equation}
    PWM_i
    \longrightarrow
    \omega_i
    \longrightarrow
    T_i,Q_i
    \longrightarrow
    IsaacSim
    \longrightarrow
    x,y,z,\phi,\theta,\psi
\end{equation}

Như vậy, controller không trực tiếp thay đổi pose của drone. Controller chỉ
apply lực vật lý, còn chuyển động thật do engine physics của Isaac Sim tích
phân.
